% ================================================================
% ==== FILE: content/m_spaces/m7.tex
% ================================================================

% ==============================
%  Ontology of Continua — Core
%  Meta-Space: M7
%  FULL MODULE — FINAL VERSION
% ==============================

\subsubsection{\texorpdfstring{$M_7$}{M_7} Übersicht}
\label{sec:m7-overview}

$M_7$ ist der Metaraum, der \emph{soziale Organisation} zulässig macht.
Während $M_6$ kognitive Kontinua unterstützt, die auf Attraktoren, Konzepten usw. basieren
protologische Operationen, $M_7$ ermöglicht:

\begin{itemize}
    \item gemeinsame Bedeutung und intersubjektive Konzepte,
    \item soziale Achsen und Koordinationsmechanismen,
    \item Kommunikationskanäle und symbolischer Austausch,
    \item Entstehung von Zyklen und Normen auf Gruppenebene,
    \item Bevölkerungspotentiale und -ströme,
    \item strukturelle Bedingungen für die Existenz sozialer Kontinua ($K_7$).
\end{itemize}

Somit ist $M_7$ die minimale Umgebung für \emph{kognitive Kopplung mehrerer Agenten}.

% ---------------------------------------------------------------
\subsubsection*{1. Zulässiger Konfigurationsraum \texorpdfstring{$\Omega(M_7)$}{\Omega(M_7)}}

\[
\Omega(M_7)=
\left\{
\big(
\mathcal{C}_i,\
A_{\mathrm{shared}},\
R_{\mathrm{comm}},\
\mathcal{N},\
C_{\mathrm{soc}},\
\Sigma_7
\Big)
\ \big| \ \text{Zulässigkeitsbedingungen gelten}
\right\}.
\]

Komponenten:

\begin{itemize}
    \item $\mathcal{C}_i$ – kognitive Kontinua von Individuen eingebettet in $M_6$,
    \item $A_{\mathrm{shared}}$ – Achsen, die gemeinsame Darstellung und Bedeutung unterstützen,
    \item $R_{\mathrm{comm}}$ — zulässige Kommunikationsbeziehungen,
    \item $\mathcal{N}$ – Netzwerkstruktur von Interaktionen,
    \item $C_{\mathrm{soc}}$ – Zyklen, die die soziale Ordnung aufrechterhalten (Normen, Rollen, Gegenseitigkeit),
    \item $\Sigma_7$ – Kohärenzbeschränkungen zur Gewährleistung der Gruppenstabilität.
\end{itemize}

Voraussetzung für die Zulassung sind:

\begin{enumerate}
    \item Existenz mindestens einer gemeinsamen semantischen Achse zwischen den Agenten,
    \item Fähigkeit, Unterscheidungen entlang $R_{\mathrm{comm}}$ zu übertragen,
    \item Stabilität sozialer Kreisläufe gegenüber Schwankungen,
    Die \item-Interaktionstopologie $\mathcal{N}$ unterstützt Kohärenz statt Fragmentierung.
\end{enumerate}

% ---------------------------------------------------------------
\subsubsection*{2. Zulässige Achsen \texorpdfstring{$A(M_7)$}{A(M_7)}}

Neue Achsen eingeführt in $M_7$:

\begin{itemize}
    \item $A_{\mathrm{shared}}$ – Achsen gemeinsamer Bedeutung / Intersubjektivität,
    \item $A_{\mathrm{comm}}$ – Kommunikationsachsen (Symbole, Signale, Normen),
    \item $A_{\mathrm{coord}}$ – Achsen koordinierter Aktion,
    \item $A_{\mathrm{role}}$ – Achsen, die sozialen Rollen entsprechen,
    \item $A_{\mathrm{norm}}$ – Achsen, die mit Normen und Regelbefolgung verbunden sind,
    \item $A_{\mathrm{coop}}$ – Kooperations-Defekt-Spektrum,
    \item $A_{\mathrm{identity}}$ – Achsen, die Gruppenidentitätsunterschiede erfassen.
\end{itemize}

Notwendige Bedingung für $K_7$:

\[
\dim(A_{\mathrm{shared}})\ge 1,\quad
A_{\mathrm{comm}}\subseteq A(M_7).
\]

% ---------------------------------------------------------------
\subsubsection*{3. Potentiale \texorpdfstring{$P(M_7)$}{P(M_7)}}

Soziale Potenziale verallgemeinern $P(M_6)$ auf eine Multi-Agenten-Organisation:

\[
P(M_7)=
(F_{\mathrm{comm}},\
F_{\mathrm{koord}},\
F_{\mathrm{Status}},\
F_{\mathrm{norm}},\
F_{\mathrm{coh}},\
F_{\mathrm{coop}}).
\]

Interpretation:

\begin{itemize}
    \item $F_{\mathrm{comm}}$ – Kommunikationspotenzial, das den Informationsfluss ermöglicht,
    \item $F_{\mathrm{coord}}$ — potenzielle stabilisierende koordinierte Aktivität,
    \item $F_{\mathrm{status}}$ — Gradienten des sozialen Einflusses und der Hierarchie,
    \item $F_{\mathrm{norm}}$ — normative Potenziale (Konformitätsdruck),
    \item $F_{\mathrm{coh}}$ — Gruppenkohärenzpotential,
    \item $F_{\mathrm{coop}}$ – energetisches Äquivalent von Kooperationsanreizen.
\end{itemize}

Ein soziales Kontinuum liegt vor, wenn:

\[
F_{\mathrm{coh}} > \Theta_{\mathrm{fragment}},
\quad F_{\mathrm{norm}} \text{ lässt stabile Zyklen zu}.
\]

% ---------------------------------------------------------------
\subsubsection*{4. Schwellenwerte \texorpdfstring{$\Theta(M_7)$}{\Theta(M_7)}}

Zu den kritischen Schwellenwerten gehören:

\begin{itemize}
    \item $\Theta_{\mathrm{comm}}$ – minimale Kommunikationsbandbreite,
    \item $\Theta_{\mathrm{shared}}$ – Schwelle für die gemeinsame Konzeptbildung,
    \item $\Theta_{\mathrm{coh}}$ – Gruppenkohärenzschwelle,
    \item $\Theta_{\mathrm{coop}}$ – Kooperationsschwelle,
    \item $\Theta_{\mathrm{norm}}$ – minimaler Druck für die Normstabilität,
    \item $\Theta_{\mathrm{fragment}}$ – Fragmentierungsschwelle bei $\partial\Omega(M_7)$.
\end{itemize}

Das Überqueren von $\Theta_{\mathrm{fragment}}$ führt zum Zusammenbruch in unabhängige $K_6$-Kontinua.

% ---------------------------------------------------------------
\subsubsection*{5. Flüsse \texorpdfstring{$J(M_7)$}{J(M_7)}}

Zulässige Flüsse:

\[
J(M_7)=
\left(
J_{\mathrm{comm}},\
J_{\mathrm{norm}},\
J_{\mathrm{coop}},\
J_{\mathrm{Identität}},\
J_{\mathrm{koord}},\
\nabla_i T_7
\right).
\]

Interpretation:

\begin{itemize}
    \item $J_{\mathrm{comm}}$ – Kommunikationsfluss über das Netzwerk,
    \item $J_{\mathrm{norm}}$ – Ausbreitung von Normen,
    \item $J_{\mathrm{coop}}$ – kooperative/Defektionsdynamik,
    \item $J_{\mathrm{identity}}$ – Flüsse, die die Gruppenidentität formen,
    \item $J_{\mathrm{coord}}$ – Flüsse, die synchronisierte Aktionen ermöglichen,
    \item $\nabla_i T_7$ – Gradient der sozialen Spannung.
\end{itemize}

Erhaltungsbeziehung (begrenzte Ressourcen an Aufmerksamkeit, Vertrauen, Engagement):

\[
\sum_i J_{\mathrm{comm},i} \le C_{\mathrm{channel}}.
\]

% ---------------------------------------------------------------
\subsubsection*{6. Zyklen \texorpdfstring{$C(M_7)$}{C(M_7)}}

Grundlegende soziale Zyklen:

\[
C_{\mathrm{soc}} =
\{
\text{Signal} \to \text{Antwort} \to \text{Feedback} \to
\text{Aktualisierung der Normen} \to \text{neues Signal}
\}.
\]

Andere Zyklen:

\begin{itemize}
    \item $C_{\mathrm{role}}$ – Rollenerwerb, Umsetzung, Verstärkung,
    \item $C_{\mathrm{coop}}$ – Kooperationsschleife (Vorteil $\to$ Vertrauen $\to$ Zusammenarbeit),
    \item $C_{\mathrm{conflict}}$ — Konflikt $\to$ Verhandlung $\to$ Lösung,
    \item $C_{\mathrm{identity}}$ — Bildung und Stabilisierung der Gruppenidentität,
    \item $C_{\mathrm{institution}}$ – Entstehung institutioneller Regeln.
\end{itemize}

Ein lebendiges soziales Kontinuum erfordert:

\[
\oint_{C_{\mathrm{soc}}} dA_{\mathrm{Zustand}} \ approx 0.
\]

% ---------------------------------------------------------------
\subsubsection*{7. Grenze \texorpdfstring{$\partial\Omega(M_7)$}{\partial\Omega(M_7)}}

Eine Konfiguration erreicht $\partial\Omega(M_7)$, wenn:

\begin{itemize}
    \item gemeinsame Achsen kollabieren (Verlust der gemeinsamen Bedeutung),
    \item-Kommunikationskanäle fallen aus oder sind überlastet,
    \item normative Zyklen brechen (keine stabilen Erwartungen),
    \item Zusammenarbeit wird nicht mehr nachhaltig,
    \item Identitätsachsen zerfallen in inkompatible Unterräume,
    \item soziale Spannung $T_7$ überschreitet Toleranzschwellen.
\end{itemize}

Das Überschreiten der Grenze führt zu einer Fragmentierung in mehrere $K_6$-Kontinua.

% ---------------------------------------------------------------
\subsubsection*{8. Beziehung zu \texorpdfstring{$M_6$}{M_6} und \texorpdfstring{$M_8$}{M_8}}

\textbf{Von $M_6$ bis $M_7$:}

Notwendige Voraussetzungen:

\[
A_{\mathrm{concept}}^{(i)} \cap A_{\mathrm{concept}}^{(j)} \neq \emptyset,
\quad
R_{\mathrm{comm}} \neq \emptyset.
\]

Für Intersubjektivität ist eine gemeinsame konzeptuelle Struktur erforderlich.

\textbf{Auf dem Weg zu $M_8$ (zivilisatorisch-technologischer Raum):}

$M_8$ stellt vor:

\begin{itemize}
    \item symbolische Systeme,
    \item Schreiben, Institutionen, Infrastruktur,
    \item technologische und wirtschaftliche Achsen.
\end{itemize}

Somit liefert $M_7$ die Grundlage der für $K_8$ erforderlichen \emph{proto-institutionellen Kohärenz}.
